\chapter{Круги и подъём на ходу}\label{ch:krugi}

\textbf{Что вы узнете:} как нарисовать круг и как лететь вверх и вперёд одновременно.

\begin{vnimanie}[Эти команды ещё не летали]
Шаги \kod{circle} и \kod{code_z} в приложении упаковываются для загрузки так же, как в LiteBee Go, но на живом дроне пока не проверены. В симуляторе они работают. В списке проверок вы увидите замечание «Шаги circle ещё не летали на живом дроне». Для первых полётов пользуйтесь командами из уроков 2--8.
\end{vnimanie}

\section{Круг: circle}

\kod{d.circle(count, cell1, cell2, seconds, rotate='cw')} — круг вокруг \emph{середины} между клетками \kod{cell1} и \kod{cell2}.
\begin{itemize}
  \item \kod{count} --- число \emph{полуоборотов}: 2 полукруга --- полный круг, дрон вернётся на место.
  \item Центр --- середина между клетками. Если обе клетки одинаковые, центр --- её центр.
  \item Радиус --- расстояние от дрона до центра. Дрон крутится вокруг центра, проходя через то место, где стоял.
  \item \kod{rotate}: \kod{'cw'} (по часовой) или \kod{'ccw'} (против).
\end{itemize}

\solfile{l15_circle.py}

Дрон стоит в клетке 14, центр круга --- центр клетки 26. Радиус 40\cm.

\begin{zametka}[Время круга]
В симуляторе число секунд --- это время всего круга. На живом дроне иначе: секунды --- время одного полуоборота, и шаг длится \kod{count} раз дольше. Круг на дроне получится дольше, чем на экране. Куда на коврике поворачивает \kod{'cw'}, в симуляторе угадано.
\end{zametka}

\section{Подъём на ходу: code\_z}

\kod{d.code_z(cell, z, seconds)} летит в клетку и одновременно меняет высоту. Число \kod{z} --- \emph{итоговая} высота, а не прирост.

\solfile{l15_climb.py}

Дрон со 60\cm поднимается до 150\cm, пока летит из клетки 13 в клетку 49, и возвращается к 60\cm по пути назад. Расстояние теперь считается по трём осям, скорость тоже.

Есть ещё \kod{circle_z} (полкруга с подъёмом), но его \textbf{нельзя загрузить в дрон}, поэтому в учебных шоу он не используется.

\begin{probuy}
  \item Тот же круг в другую сторону (\kod{'ccw'}).
  \item Половина круга (\kod{count} = 1): дрон оказывается напротив своего места. Где именно, если он стартует в клетке 14 с центром 26? Посадите его через \kod{land_on}.
  \item «Горка»: со 60\cm до 140\cm по пути из 13 в 16 и обратно.
  \item Два дрона в клетках 14 и 38 крутятся вокруг клетки 26, один напротив другого. Почему они не столкнутся?
\end{probuy}
\otvety{15}
